\subsubsection{Adaptive Time Stepping}
\label{SEC:Time:Algorithms:Adaptive}

A second set of ERK methods exist within GFR which belong to the family of
embedded-pair explicit Runge-Kutta (EPERK) methods.
%
EPERK methods contain an additional Runge-Kutta method of lower order than the
primary method that is simultaneously evaluated with minimal additional cost,
thus an embedded pair of methods.
%
At the end of a given time step, the difference between these two methods can be
viewed as the error in the primary time-stepping method.
%
This error can be used as the input to a PID controller to automate the
selection of an optimal time step size for the following
iteration.\cite{Berzins_1995,Ware_1995,Ranocha_2022}

%%%%%%%%%%%%%%%%%%%%%%%%%%%%%%%%%%%%%%%%%%%%%%%%%%%%%%%%%%%%%%%%%%%%%%%%%%%%%%%%

The implementation and usage of the optimized EPERK schemes follows the work of
\linecite{Ranocha_2022}.
%
The EPERK methods use the low-storage 3S*$\!{}_+$ algorithm shown in
\cref{ALG:Time:Algorithms:EPERK} that requires the same three registers as the
ERK 3S* low-storage algorithm plus a fourth register, $S_4$, for the error
estimate.
%
%\begin{algorithm}[tb]
      \centering
    \caption{Embedded-pair low storage implementation (3S*$\!{}_+$)}
    \label{ALG:Time:Algorithms:EPERK}
    \begin{algorithmic}[1]
      \STATE $S_1 \leftarrow u^n ;
              S_2 \leftarrow 0 ;
              S_3 \leftarrow u^n ;
              S_4 \leftarrow u^n $
      \FOR{$i=1:s$}
      \STATE $\tau \leftarrow t^n + c_i \Delta t$
      \STATE $S_2 \leftarrow S_2 + \delta_i S_1$
      \STATE $R_i \leftarrow \bm{\mathrm{R}}\left(\tau,S_1\right)$
      \STATE $S_1 \leftarrow \gamma_{1,i} S_1 + 
                             \gamma_{2,i} S_2 +
                             \gamma_{3,i} S_3 + 
                      \beta_i \Delta t R_i$
      \STATE $S_4 \leftarrow S_4 + 
                      \hat{b}_i \Delta t R_i$
      \ENDFOR
      \STATE $u^{n+1} \leftarrow S_1$
      \STATE $\hat{u}^{n+1} \leftarrow S_4$
      \IF{$\hat{b}_{s+1} \ne 0$}
      \STATE $\hat{u}^{n+1} \leftarrow \hat{u}^{n+1} +
            \hat{b}_{s+1} \Delta t \bm{\mathrm{R}}\left(t^{n+1},u^{n+1}\right)$
      \ENDIF
    \end{algorithmic}

\end{algorithm}
%
%
The quantity $\hat{u}^{n+1}$ represents the solution as updated by the embedded
lower-order scheme using the additional coefficients $\hat{b}_i$.
%
The time stepping error from the embedded RK pair is given by
%
\begin{equation*}
  \varepsilon_{n+1} = \frac{1}{w_{n+1}},\quad
  w_{n+1} = \frac{1}{m} \sum_{i=1}^m \left( 
     \frac{u_i^{n+1} - \hat{u}_i^{n+1}}
          {\mathtt{atol}+\mathtt{rtol}\!\cdot\!
           \max\left\{\left|u_i^{n+1}\right|,
                      \left|\hat{u}_i^{n+1}\right|\right\}}
          \right)
 %\label{eq:eperk:error}
\end{equation*}
%
where $m$ is the total degrees-of-freedom, and $\mathtt{atol}$ and
$\mathtt{rtol}$ are the respective absolute and relative error tolerances.
%
Finally, the time stepping error is used with a PID controller to determine the
time step size for the next iteration
%
\begin{equation*}
  \Delta t^{n+1} = \varepsilon^{\beta_{1} / k}_{n+1}
                   \varepsilon^{\beta_{2} / k}_{n}
                   \varepsilon^{\beta_{3} / k}_{n-1} \Delta t^n
 %\label{eq:eperk:controller}
\end{equation*}
%
where $k$ is the order of the EPERK method, and $\beta_{i}$ are the controller
parameters.
%
The error tolerances and controller parameters used in this work are those
recommended in \linecite{Ranocha_2022}.

%%%%%%%%%%%%%%%%%%%%%%%%%%%%%%%%%%%%%%%%%%%%%%%%%%%%%%%%%%%%%%%%%%%%%%%%%%%%%%%%

%Each of these optimized methods contains two versions, a base version and a
%first-same-as-last (FSAL) version.
%%
%In FSAL methods, the last evaluation of the right-hand-side function needed by
%the embedded scheme is equivalent to the right-hand-side function needed for the
%first stage of the following time step.
%%
%Therefore, this redundancy in FSAL methods can be exploited to reduce the number
%of right-hand-side evaluations per time step by one.

%%%%%%%%%%%%%%%%%%%%%%%%%%%%%%%%%%%%%%%%%%%%%%%%%%%%%%%%%%%%%%%%%%%%%%%%%%%%%%%%

\NOTE Adaptive time stepping cannot be used with limiting\slash filtering.
%
The limiting\slash filtering step, which modifies the solution variables, occurs
immediately after the completion of a Runge-Kutta stage.
%
This means the right-hand-side that was used by the EPERK method to update the
solution variables is no longer consistent with the limited\slash filtered
solution variables.
%
This inconsistency also applies to the embedded ERK method which ultimately
causes the adaptive time stepping algorithm to fail when using limiting\slash
filtering.
%
Being a global time stepping method and its incompatibility with limiting\slash
filtering, this option is suitable for time-accurate simulations of low-speed
flows.
%
Further details about EPERK methods and adaptive time stepping can be found in
\linecite{Spiegel_2025}.

%%%%%%%%%%%%%%%%%%%%%%%%%%%%%%%%%%%%%%%%%%%%%%%%%%%%%%%%%%%%%%%%%%%%%%%%%%%%%%%%

As with the CFL-based global time stepping, a unique time step size is computed
for each cell, and the minimum time step size across all grid cells is taken to
be the global time step size that all cells use to update their solutions in
time.
%
When it works as intended, this method is the easiest to use because the
adaptive algorithm produces a near-optimal time step size with zero input from
the user.
%
However, testing has found that it tends to not work more than it works as
intended.

%%%%%%%%%%%%%%%%%%%%%%%%%%%%%%%%%%%%%%%%%%%%%%%%%%%%%%%%%%%%%%%%%%%%%%%%%%%%%%%%

